%% kapa: exam=1 round=28 year=2017 period=1 form=A subject=economics subject_ko=경제학원론 seq=7 number=47
%% source: Q-Net 감정평가사 기출문제 (cst00309) artlSeq=5199603 file=extracted/1008/5199603/2017년 제28회 감정평가사 1차시험 문제지/1교시 A형.pdf
%% conversion: kordoc 4.24.51 -> tools/gen1.py (auto); review=proofread-sonnet; flags=none
\begin{question}{47}
주어진 소득으로 X재, Y재 두 재화만을 소비하는 甲의 효용함수가 $U=x^{1/3}y^{2/3}$일 때, 설명으로 옳지 않은 것은? (단, $x$는 X재 소비량, $y$는 Y재 소비량이며, 소득과 두 재화의 가격은 0보다 크다.)
\begin{choices}
  \item X재는 정상재이다.
  \item Y재는 정상재이다.
  \item 甲의 무차별곡선은 원점에 대해 볼록하다.
  \item 두 재화의 가격비율에 따라 어느 한 재화만 소비하는 결정이 甲에게 최적이다.
  \item 두 재화의 가격이 동일하다면 Y재를 X재보다 많이 소비하는 것이 항상 甲에게 최적이다.
\end{choices}
\end{question}
\answer{4}
